\section{Pair and triple scales}
\label{p2:scales}

The first deviation from Poisson is determined by rational coincidences
between two separated windows. We keep the complete finite coefficient:
its lower-order terms matter before the next channel becomes visible.
The counting estimates behind this calculation are proved in
Sections~\ref{p1:arith:appendix} and~\ref{p1:arith:section}.
For the pair profiles write $\rho(I)=|I|L-\dim\sum_{x\in I}S_x$
and let $\eta(I)$ indicate affine compatibility, so that
$q(I)=p^{|I|}\eta(I)2^{\rho(I)}$. This raw nullity includes local
kernels and is distinct from the effective relative nullity used in
Section~\ref{p13:exact:section}.

\subsection{The two leading rational channels}

An exact unit of the reduced channel $a:b$, $a<b$, is a pair of values
$(at,bt)$ in the two windows.  A relative rational relation selects an even
number of these units.  If the first and last units are $t$ and $t+d$, the
number of selections is
\[
 A_1=1,\qquad A_d=2^{d-2}\quad(d\geq2).
\]
Summing over positions in one window gives $L+1-D$ for a selection of
diameter $D$.  Retaining the predecessor sign changes this to
\begin{equation}\label{p1:arith:diametersign}
 g_L(D)=(L-1-D)\ind_{\{0\leq D\leq L\}}.
\end{equation}
Indeed exactly one of the $L+1-D$ placements has its first selected vertex
at the predecessor; its sign changes from $+1$ to $-1$.  The value
$g_L(L)=-1$ is part of the formula.

Define
\begin{align*}
 U_{a,b}(L)&=\sum_{d=1}^{\lfloor L/b\rfloor}
            A_d(L+1-ad)(L+1-bd),\\
 S_{a,b}(L)&=\sum_{d=1}^{\lfloor L/b\rfloor}
            A_d(L-1-ad)(L-1-bd).
\end{align*}
For $L=bm+r$, $0\leq r<b$ and $m\geq1$, both sums are evaluated by
\begin{equation}\label{p1:arith:channelidentity}
 \begin{split}
 &\sum_{d=1}^{m}A_d(L+s-ad)(L+s-bd)\\
 &\quad=2^{m-1}\big[(b-a)(b+r+s)m+(r+s)^2
                   +(a+b)(r+s)+3ab\big]\\
 &\qquad\quad -(a+b)(L+s)-ab,\qquad s\in\{-1,1\}.
 \end{split}
\end{equation}
When $m=0$ the defining sums are empty and equal to zero.
To verify the identity for $m\geq1$, expand the product and use
\[
 \sum A_d=2^{m-1},\quad
 \sum dA_d=(m-1)2^{m-1}+1,\quad
 \sum d^2A_d=(m^2-2m+3)2^{m-1}-1.
\]
In particular, for $m\geq1$,
\begin{align}
 S_{2,3}(L)&=2^{m-1}\big[(r+2)m+r^2+3r+14\big]-5L-1,
       &&L=3m+r,\quad0\leq r<3,\label{p1:arith:S23}\\
 S_{3,4}(L)&=2^{m-1}\big[(r+3)m+r^2+5r+30\big]-7L-5,
       &&L=4m+r,\quad0\leq r<4.\label{p1:arith:S34}
\end{align}

\begin{theorem}[Complete pair profile]\label{p1:arith:pairs}
The sums below are over ordered separated pairs in $I_N$:
\begin{align}
 \sum_{|x-y|>L}(2^{\rho(x,y)}-1)
 &=\frac N3\{U_{2,3}(L)+U_{3,4}(L)\}
      +O_{C,\varepsilon}(N^{6/5+\varepsilon}),\label{p1:arith:positiveprofile}\\
 \sum_{|x-y|>L}(\eta(x,y)2^{\rho(x,y)}-1)
 &=\frac N3\{S_{2,3}(L)+S_{3,4}(L)\}
      +O_{C,\varepsilon}(N^{6/5+\varepsilon}).\label{p1:arith:signedprofile}
\end{align}
Writing $\mu_N=\E Z_N$ and $\alpha_N=\operatorname{Var}(Z_N)-\mu_N$, one has
\begin{equation}\label{p1:arith:cumulantexpansion}
 \alpha_N=\frac{\lambda_N^2}{3N}\{S_{2,3}(L)+S_{3,4}(L)\}
            +O_{C,\varepsilon}(N^{-4/5+\varepsilon}).
\end{equation}
Thus, with $r_3=L\bmod3$,
\begin{equation}\label{p1:arith:leadingphase}
 \alpha_N\sim
 \frac{r_3+2}{18\,2^{r_3/3}}\lambda_N^{5/3}L N^{-2/3}.
\end{equation}
After the entire first-channel term is removed, the next term is
\begin{equation}\label{p1:arith:secondphase}
 \alpha_N-\frac{\lambda_N^2}{3N}S_{2,3}(L)
 \sim\frac{r_4+3}{24\,2^{r_4/4}}\lambda_N^{7/4}L N^{-3/4},
 \qquad r_4=L\bmod4.
\end{equation}
\end{theorem}

The exact-unit calculation gives the displayed channel sums. Both
orientations of each of the channels $2:3$ and $3:4$ have a translation
interval of total length $N/3+O(B)$. Private prime coordinates identify
their rational kernels; the exceptional normalizations and all remaining
channels have the stated residual mass. Section~\ref{p14:comp:complete-profile}
gives the complete proof, including the endpoint, contact, and
mean-transfer estimates.

The subtraction in \eqref{p1:arith:secondphase} must retain the whole finite
expression $S_{2,3}(L)$.  Its own term of size $N^{-2/3}$ without the factor
$L$ is larger than $LN^{-3/4}$.  The two exact channel polynomials together
depend on $L\bmod12$.


\subsection{A computable two-channel reference}

An entirely explicit reference has a different optimal rate.
Retain the complete \(2:3\) and \(3:4\) channel sums, and set
\begin{equation}
 \widehat\alpha_N=\frac{\lambda_N^2}{3N}
             \{S_{2,3}(L)+S_{3,4}(L)\}.
 \label{p1:eq:alpha-explicit}
\end{equation}
The sums have been defined above and are evaluated by
\eqref{p1:arith:S23}--\eqref{p1:arith:S34}.

\begin{corollary}[An explicit probability reference]
\label{p1:cor:explicit}
For every \(\varepsilon>0\), uniformly in the critical band,
\begin{equation}
 \left\|\Law(Z_N)-\Pi_{\lambda_N}
       -\frac{\widehat\alpha_N}{2}
          \Pi_{\lambda_N}*\Delta^{*2}\right\|_1
       \ll_{C,\varepsilon}N^{-4/5+\varepsilon}.
 \label{p1:eq:explicit-law}
\end{equation}
For all sufficiently large \(N\), the reference
\(Q_{\lambda_N,\widehat\alpha_N}\) is a probability measure.
Its mass at \(k\ge0\) is
\begin{equation}
 e^{-\lambda_N}\frac{\lambda_N^k}{k!}
 \left[
  1+\frac{\widehat\alpha_N}{2}
       \left\{\frac{k(k-1)}{\lambda_N^2}
                    -\frac{2k}{\lambda_N}+1\right\}
 \right].
 \label{p1:eq:explicit-masses}
\end{equation}
All parameters in this reference are determined by \(N,L\).
\end{corollary}

\begin{proof}
Theorem~\ref{p1:arith:pairs} gives
\(\alpha_N-\widehat\alpha_N=O_{C,\varepsilon}(N^{-4/5+\varepsilon})\).
Since \(\|\Pi_{\lambda_N}*\Delta^{*2}\|_1\le4\), replacing the
coefficient costs at most \(2|\alpha_N-\widehat\alpha_N|\).
Apply Theorem~\ref{p1:thm:main}.
The replacement of \(\mu_N\) by \(\lambda_N\) costs
\(O(N^{-1+\varepsilon})\), which is absorbed here.
The finite formulas give
\(\widehat\alpha_N\sim\alpha_N>0\) and
\(\widehat\alpha_N\to0\), uniformly in the band.
The positivity criterion in Lemma~\ref{p1:close:positivity}
therefore applies, since \(\lambda_N\) remains in a compact
subinterval of \((0,\infty)\).
\end{proof}

The exponent \(4/5\) in \eqref{p1:eq:explicit-law} is optimal
for this two-channel reference. Its leading error is half the
height-five contribution times
\(\|\Pi_{\lambda_N}*\Delta^{*2}\|_1\), as proved in
Corollary~\ref{p13:channels:height-five}.
The explicit parameters do not supply a numerical value of
the constant in the remainder or a numerical threshold for \(N\).
The same two-channel coefficient also gives a positive Hermite
reference; see Corollary~\ref{p1:cor:compound}.

\paragraph{The probability of no retained start.}
Evaluating \eqref{p1:eq:explicit-law} at zero gives the immediate
consequence
\begin{equation}
 \P(Z_N=0)=e^{-\lambda_N}
       \left(1+\frac{\widehat\alpha_N}{2}\right)
       +O_{C,\varepsilon}(N^{-4/5+\varepsilon}),
 \label{p1:eq:no-start}
\end{equation}
because \((\Pi_{\lambda_N}*\Delta^{*2})\{0\}=e^{-\lambda_N}\).
The positive correction shows an effect of the arithmetic pairs
that is absent from the mean alone. Here \(Z_N=0\) means that
no retained run start lies in \(I_N\). A run starting before
\(N\) may still intersect the interval, and a run starting in
\(I_N\) may continue beyond \(2N\).
The exact-parameter refinement is Corollary~\ref{p13:optimal:zero}.

\subsection{Why triple interactions are smaller}

The third factorial cumulant is
\begin{equation}
 \kappa_{3,N}=\E(Z_N)_3-3\mu_N\E(Z_N)_2+2\mu_N^3.
 \label{p13:third:definition}
\end{equation}
It is distinct from the ordinary third cumulant, which equals the mean
for a Poisson variable.

\begin{theorem}[Optimal power of the third factorial cumulant]
\label{p13:third:main}
For each fixed critical band $C>0$, there is a constant $c_C>0$ such
that, for all sufficiently large $N$,
\begin{equation}
 c_C N^{-41/30}\leq\kappa_{3,N}
              \ll_{C,\varepsilon}N^{-41/30+\varepsilon}
              \qquad(\varepsilon>0).
 \label{p13:third:two-sided}
\end{equation}
Both bounds are unconditional.  Thus $41/30$ is the optimal power
exponent for this factorial cumulant.  The theorem does not assert an
asymptotic equivalent or identify a logarithmic factor for the full
cumulant.
\end{theorem}

The complete proof is given in Section~\ref{p13:third:section}.
It bounds all admissible triples, with the exact marginals retained.
The following form is the input to the reconstruction.

\begin{corollary}[Absolute third-order activity]
\label{p13:third:activity}
For every $\varepsilon>0$,
\begin{equation}
 A_3:=\sum_{|T|=3}|c(T)|
            \ll_{C,\varepsilon}N^{-41/30+\varepsilon}.
 \label{p13:third:absolute}
\end{equation}
\end{corollary}

\begin{proof}
Equation~\eqref{p13:activity:kappa-comparison} gives
$A_3=K_3^++O(N^{-3/2+\varepsilon'})$.
Use the upper bound of Theorem~\ref{p13:third:main} and reduce the internal margin.
\end{proof}


The geometric step in the upper bound changes the scale at which points
are counted. Write $d_{xy}=\dim(S_x\cap S_y)$ and suppose
$d_{xy}>((2/3)\log2+\delta)B$ for a fixed $\delta>0$.
Short-relation extraction yields a space of dimension $\gg_\delta B$
and a common curve through $(x,y)$, of degree
$D\le B^{a+o(1)}$ for some fixed $a<1/2$.
For the relations in a subspace of codimension at most one, the
associated products are squares in the curve's function field.
At a selected noncritical offset $u$, the simple zeros of $X+u$
must therefore be matched by zeros of factors $Y+v$ from the
other window. After $O(D^2)$ exceptional offsets
are removed, these fibres give $\gg_\delta B$ integer points in
$[-B,B]^2$. Their coordinates are logarithmic-size offsets, although
the original centres have size $N$. For $D\ge2$, the degree-uniform
point bound gives only $o(B)$ points, a contradiction. For $D=1$,
matching the roots gives a small rational slope and an alignment.
Thus nonaligned pairs obey the rank bound of
Theorem~\ref{p13:rig:nonaligned}; the full argument is in
Section~\ref{p13:rig:section}.

The configuration $10t,12t,15t$ explains the third-order power.
Multiplication by $6,5,4$ gives a common center $60t$ and offset
steps $6,5,4$. The canonical relative nullity is
\[
 \rho_L=\left\lfloor\frac{L-1}{6}\right\rfloor+
 \left\lfloor\frac{L-1}{3}\right\rfloor+
 \left\lfloor\frac{L-1}{5}\right\rfloor-
 \left\lfloor\frac{L-1}{15}\right\rfloor
 =\frac{19}{30}L+O(1).
\]
A compatible triple with no additional square relations therefore has
weight $N^{-71/30+o(1)}$. For a concrete phase, take $L=30m$ and $N=2^L$.
Then
\[
 \begin{aligned}
 \rho_L&=(5m-1)+(10m-1)+(6m-1)-(2m-1)=19m-2,\\
 q(\{10t,12t,15t\})&=p^3 2^{\rho_L}=\frac14N^{-71/30}.
 \end{aligned}
\]
There are $N/30+O(1)$ translations in the dyadic interval.
Proposition~\ref{p13:third:lower} shows that all but $N^{1/2+o(1)}$
of them have private prime coordinates and compatible predecessor
signs. On this population, the largest pair exponent subtracted in
the centered coefficient is $10m-1$, so $c(T)=q(T)(1+o(1))$.
This family thus contributes $(1/120+o(1))N^{-41/30}$ to $K_3^+$;
the coefficient is specific to this family, not to the full cumulant.
The upper bound also needs pair rigidity and an exhaustive partition
of the other triples. It is not a consequence of this example alone.

The distinction between pairs and triples is decisive:
$\alpha_N^2=N^{-4/3+o(1)}$, whereas
$A_3\ll N^{-41/30+\varepsilon}$ and $41/30-4/3=1/30$.
The reconstruction will retain the first pair term and allow its
square in the error. The fourth-order test will prove that this
remaining scale cannot be improved in power for the quadratic reference.
